\section{Funcția hash}

Funcția hash propusă în această lucrare se bazează pe toate conceptele prezentate în secțiunile anterioare.
Structura Sponge este aplicată pe blocuri, iar apoi prin intermediul arborelui Merkle se obține hash-ul întregului mesaj.
Dimensiunea stării interne este de 1024 biți, împărțită în 512 biți pentru rată și 512 biți de capacitate.
Transformarea internă a stării este realizată printr-o secțiune haotică, completată de metode clasice de confuzie și difuzie. 
Pentru fragmentele următoare de pseudocod, considerăm următoarele convenții de notație:
\begin{itemize}
    \item $S$ reprezintă starea internă, formată din 32 de cuvinte de 32 de biți, $S_i$ fiind cuvântul de pe poziția $i$.
    \item $a \oplus b$ reprezintă XOR-are (disjuncție exclusivă) bit cu bit.
    \item $\rotl{a}{k}$ reprezintă rotația circulară spre stânga a cuvântului $a$ cu $k$ poziții.
    \item $a \bmod b$ reprezintă restul împărțirii lui $a$ la $b$.
    \item $l$ reprezintă lungimea hash-ului final, exprimată în biți.
\end{itemize}

\subsection{Preprocesarea}
Secțiunea de preprocesare are rolul de a pregăti mesajul inițial pentru absorbția în construcția Sponge.
Astfel, lungimea mesajului preprocesat trebuie să fie un multiplu al ratei de 512 biți.
Considerând că mesajul inițial este $M$, alcătuit din $M=M_1 \| M_2 \| \ldots \| M_n$, unde $M_i$ reprezintă un byte, mesajul preprocesat va fi:

\begin{equation}
    M' = M \| \texttt{0x80} \| \texttt{0x00} \| \ldots \| \texttt{0x00} \| \texttt{0x01},
\end{equation}

\subsection{Generarea frunzelor}
După preprocesare, mesajul este împărțit în blocuri de dimensiune configurabilă.
Fiecare bloc $B_i$ va reprezenta o frunză în arborele Merkle, iar hash-ul său va fi calculat prin aplicarea construcției Sponge.
Trebuie menționat faptul că dimensiunea unui bloc este un parametru configurabil.

\subsection{Construcția Sponge}
Construcția Sponge este aplicată pe fiecare bloc $B_i$ pentru a genera hash-ul.
Fiecare bloc de 512 biți este XOR-at cu partea de rată a stării interne, după care se aplică permutarea haotică.
Permutarea $f$ a construcției Sponge este aplicată după fiecare absorbție a unui bloc.
Starea internă este împărțită în 16 secțiuni, fiecare formată din două cuvinte de 32 de biți, unul din partea de rată și celălalt din partea de capacitate.
Ulterior, cele două sisteme haotice sunt aplicate secvențial pe fiecare secțiune de 20 de ori.
După ce toate secțiunle au fost procesate, o etapă de difuzie amestecă secțiunile între ele.
Astfel, mai întâi se amestecă cuvintele vecine din fiecare jumătate, apoi se amestecă cele două jumătăți între ele.
Amestecarea se realizează prin rotații de biți și XOR-are, pentru asigura propagarea schimbărilor în întreaga stare internă.

\begin{algorithm}[!hbpt]
\caption{Permutarea stării}
\label{alg:chaos}
\KwData{starea internă $S_i,\ i \in \overline{0,15}$}
\KwResult{starea nou permutată $S'_i$}
\For{$i \gets 0$ \KwTo $15$}{
    $x \gets S_i$\;
    $y \gets S_{i+16}$\;
    \For{$r \gets 1$ \KwTo $20$}{
        \tcp{harta brutarului}
        \eIf{$x < 2^{31}$}{
            $x \gets x \ll 1$\;
            $y \gets y \gg 1$\;
        }{
            $x \gets (\lnot x) \ll 1$\;
            $y \gets ((\lnot y) \gg 1) \lor 2^{31}$\;
        }
        \tcp{harta omului din turtă dulce}
        \eIf{$x \geq 2^{31}$}{
            $a \gets -x$\;
        }{
            $a \gets x$\;
        }
        $(x,\ y) \gets (a - y,\ x)$\;
    }
    $S_i \gets x$\;
    $S_{i+16} \gets y$\;
    \tcp{cuplarea coloanelor: reinjectare în capacitatea următoare}
    $t \gets (i + 1) \bmod 16 + 16$\;
    $S_t \gets S_t \oplus (\rotl{x}{9})$\;
    $S_t \gets S_t \oplus (\rotl{y}{17})$\;
}
\tcp{difuzie: pas înainte și pas înapoi în fiecare jumătate}
\ForEach{$b \in \{0, 16\}$}{
    \For{$i \gets 0$ \KwTo $15$}{
        $S_{b+i} \gets S_{b+i} \oplus (\rotl{S_{b + (i+1) \bmod 16}}{13})$\;
    }
}
\ForEach{$b \in \{0, 16\}$}{
    \For{$i \gets 15$ \KwTo $0$}{
        $S_{b+i} \gets S_{b+i} \oplus (\rotl{S_{b + (i+15) \bmod 16}}{11})$\;
    }
}
\tcp{difuzie: pas de amestecare între rată și capacitate}
\For{$i \gets 0$ \KwTo $15$}{
    $S_i \gets S_i \oplus (\rotl{S_{i+16}}{7})$\;
}
\For{$i \gets 0$ \KwTo $15$}{
    $S_{i+16} \gets S_{i+16} \oplus (\rotl{S_i}{19})$\;
}
\end{algorithm}

\subsection{Construirea arborelui Merkle}
Construcția Sponge și algoritmul descris anterior sunt utilizate în cadrul arborelui Merkle.
Dacă nodul curent este o frunză, atunci datele absorbite sunt blocurile $B_i$ ale mesajului preprocesat $M'$.
Dacă nodul curent este un nod intern, atunci datele absorbite sunt concatenarea ratelor de la toții copii nodului curent.
Trebuie menționat faptul că numărul maxim de copii ai unui nod intern este configurabil.
La finalul procesului, rădăcina arborelui Merkle conține hash-ul întregului mesaj.
Înainte de faza de stoarcere a hash-ului din construcția Sponge, se aplică un număr de runde intermediare, unde nu este absorbit niciun bloc de date.
Numărul de runde intermediare $n$ este determinat de lungimea mesajului inițial și de dimensiunea finală a hash-ului prin următoarea formulă:

\begin{equation}
    n = 9 + (|M| \bmod 37 + |l| \bmod 31) \bmod 11
\end{equation}

La final, hash-ul este extras din partea de rată a stării interne.
Se execută faza de stoarcere până când se obțin suficienți biți pentru a forma hash-ul.
Având în vedere că rata este de 512 biți, dacă lungimea hash-ului final $l$ nu este multiplu de 512, biții rămași sunt trunchiați.